\section{Unchecked inputs and remaining gates}
\label{sec:remaining}

The following entries are deliberately outside the checked-results
part. They are not assigned a proof status by analogy with a nearby
proved result.

\paragraph{U01: the requested Hooley-based proof.}
A02 proves the arithmetic consequence of M. A03 proves the exact
prime-degree geometric implication, and A11 weakens its sufficient
condition to an odd prime power $n$ with
$\langle2,-1\rangle=(\mathbb Z/n)^\times$ and an exact-order-$n$
theta point. A15 and A18 weaken this further: in characteristic two and every
genus at least two, a cyclic cover of any odd degree with primitive
factor of geometric $2$-rank zero suffices. In genus two this is
equivalent to theta meeting every primitive doubling orbit of the
defining cyclic torsion subgroup. A18 removes the exceptional
elliptic double quotient by an exact Deuring--Shafarevich and
conductor-rank calculation.
A04 refutes the inference from
two independent density statements, and A06 refutes the universal
joint primitive-root theta assertion itself. A05 proves A1 by a
different route. No claim is made here that a proof of A1 through
Hooley has been found. A strategy restricted to positive-two-rank
genus-two curves still requires a curve-dependent incidence theorem.
A07 refutes the positive-rank version over finitely generated fields
with transcendental moduli, not over finite fields. A08 shows that
no fixed finite prime list works over all finite fields; in the
ordinary case it excludes any fixed finite list of odd degrees.
A09 supplies an exact degree-$13$ rank-one witness. A13 supplies an
ordinary example of prime degree $4159$, with $2$ of index two in
the unit group. A16 supplies a rank-one example of degree $15$
whose full new factor has positive $2$-rank but whose primitive
factor has rank zero. A10 excludes order-three theta incidence on every nonordinary
genus-two base. None of
these results interchanges the finite-list and all-primes quantifiers;
in higher genus it requires nilpotence of every relevant Frobenius
orbit composite, not merely one nonordinary arrow.
A18 now permits any odd conductor with geometric primitive rank zero.
A19 makes all exponents for each specified finite odd prime support
an exact finite search. A21--A22 show that neither generic
positive rank nor any fixed finite support forces success, even
when mixed orders are allowed. A23 refutes universal success of
the older large-prime quotient inequality over constant extensions.
The remaining existence assertion must therefore allow the prime
support to depend on the particular finite-field curve, or use a
separating construction outside these sufficient tests.

A26 weakens the separate rational quotient criterion using even
augmentation. A27 nevertheless gives a rank-one curve over $\mathbf F_4$
where its numerical hypothesis fails for every extension and every
primitive-two prime. A28 recalls the recovered canonical-class criterion,
and A29 strictly improves its prime-level detection: a nonzero
$r$-primary component of $[K_C-2P]$ suffices if
$\operatorname{ord}_r(p)$ is even, with a cyclic $r^a$ cover and a
prime-to-$p$ rational top Prym. This is unconditional and needs no
prime-distribution theorem. The actual-curve strictness example in A29
shows why passage to prime powers matters even when all prime-level
canonical detectors fail. The residual points have canonical class
supported on the characteristic prime and odd primes of odd
multiplicative order, with the prime-two issue separate in odd
characteristic. In characteristic two this includes every Weierstrass
point and the purely two-primary non-Weierstrass examples of A09.
Constant extensions preserve that primary support. Other covers may
still work; no absence conclusion follows from this residual condition.

A30 refutes a proposed direct higher-genus extension of the
``every detecting quotient is good'' conclusion: a genus-three
rank-one curve has a detecting constant cyclic cubic quotient whose
relative class number is even. It does not exclude other good
quotients. The proposed unramified double-cover repair is unavailable
to the geometric prime-to-two fundamental group and does not preserve
the descending detector. A31 shows that the rank-one curve of A27
nevertheless has an exact signed-doubling theta witness of order seven.
All six of its original rational points permanently miss A29's
canonical detector, yet the geometric witness can be pointed after
a finite extension. Thus the sufficient tests have distinct limits;
their failure is not a nonexistence result.

The approaches actually tested were: fixed-base Hooley density;
Goldfeld's average over bases; Lenstra's fixed finite Chebotarev
constraints and positivity; prime-to-characteristic torsion orders
from Pop--Sa\"idi; a single singular Frobenius arrow in higher genus;
scalar-Frobenius packet counts; generic ordinary and fixed-rank
Prym theorems with finite-type specialization; a direct degree-$13$
theta certificate; prime-power character blocks; primary torsion
projections on theta; the Lang quotient and Weil bound; growth of
fixed-prime primary torsion; the composite primitive projector;
an exact order-$15$ search and certificate; the Hooley--Lenstra
family $n=3q$; the elliptic-exception rank calculation; finite
prime-support torsion via Frobenius translation orbits; the
all-three-primary obstruction; constant twists with an unspecified
split base point; the parity augmentation test; its all-extension rank-one obstruction;
inverse-character canonical detection and its prime-power refinement;
and the independent cuspidal route.
Their respective endpoints are recorded in A02--A31 and the
retained B.28--B.29 pointing audit. Poonen's
projection theorem supplies an order $\ell^r m$, with uncontrolled
prime factors of $m$; multiplying by $m$ moves theta, and passing to
an intermediate cover leaves proper-conductor character blocks.
Those are failures of those two strategies. A15 bypasses the
proper-conductor factors by a product projector on its stated
geometric locus, and A18 removes its geometric exception in
characteristic two. What remains is the same-subgroup orbit-coverage
condition, or the stronger single-point unit-group condition.
A17 gives density $\mathcal A/2$ to the candidate primes
$q\equiv5\pmod8$ for $n=3q$, conditionally on M, without proving
theta incidence for any particular curve. A12 gives an unconditional
curve-dependent primary-torsion test; A14 proves that enlarging the
constant field at one fixed prime cannot eventually make that
particular sufficient test succeed. Neither statement is a theorem
of nonexistence of theta torsion. The recovered Successor A.10
rational prime-cover criterion is separately distinguished in A14.
A19 now makes all primary exponents a finite exact search for
each prescribed finite odd prime support. A20 proves that even
all powers of three fail on the rank-one curve of A09, which
nevertheless has its order-thirteen witness. The prime supports
remain unbounded; no theorem proves that one of these finite
searches succeeds for every positive-rank curve.
The B.28--B.29 rational-twist alternative supplies a good split
point without prescribing it. Its good-point density estimate
does not force such a point into an arbitrarily prescribed fibre,
and extending a good twist can destroy its rational torsion
condition. These failures are recorded separately from geometric
primitive-rank-zero incidence.
The cuspidal route
settles point separation, while the scalar route refutes the
universal auxiliary existence assertion. None of the other routes
silently fills the curve-dependent incidence gap.

\paragraph{U02: the literal primitive-cover problem.}
ANABELIAN\_A1 B.24 allows the degree-one cover with convention
$P_1=J_C$. The B.23 counterexample has $p$-rank zero and is therefore
covered by that convention. The problem for positive $p$-rank
remains open in the recovered run. Similarly ANABELIAN\_A1 B.5 and
B.32 ask for prescribed-point rational class-number witnesses. A29
now settles the stated nonzero canonical-primary component cases,
while their universal existence remains open.
A bound for a prime in a number-field Frobenius class does not
produce the required covering geometry or local unit conditions.
No assertion is made that these separate questions follow from A05.

\paragraph{U03: complete source coverage.}
The two missing registers account for 390 full statements. Only
47 distinct A1-successor IDs and 50 distinct GT-tower IDs are visible
in the supplied summaries and flag excerpts. The missing LM7
proof files account for another 209 proof blocks. Exact-title,
variant-title and content searches recovered four relevant archives,
but not the two missing sources. A final prior-context recovery
attempt was unavailable because that retrieval service was disabled.
The next concrete input is the complete
\texttt{ANABELIAN\_A1\_SUCCESSOR} register/paper and the complete
\texttt{GT\_TOWER\_CHECK} register/paper, together with the omitted
LM7 proof sections. Their absence is an evidentiary gate, not a
mathematical theorem beyond GRH.

\paragraph{U04: infinite towers and unspecified finite data.}
For a recovered row with finite affine monodromy, G01 gives a
prime bound after the actual finite image, cocycle, precision,
and excluded primes are fixed. It does not bound the precision
needed to separate an infinite object, the height of a Hilbert
specialization, or a Baire-category construction. A finite test of
full adic semisimplifications is available under the hypotheses of
G03, with exact computable traces as an additional algorithmic input.
Its applicability to an unspecified row must be checked separately.

\paragraph{U05: scope of the effective open-image bounds.}
The former descent gate for LM8 A.32--A.33 is closed by their actual
models $F_{\pm6}/\mathbb Q$ and the checked refinement G06.
The conditional Mayle--Wang bound applies over $\mathbb Q$; after
restriction to the proof's number field $L$, G06 proves an index
bound $[L:\mathbb Q]$, not necessarily full image.
G-LOMBARDO is an unconditional explicit adelic index bound over
arbitrary number fields. The hypotheses must still be checked for
any other factor; CM factors have CM monodromy rather than full
$\operatorname{GL}_2$ monodromy. No corresponding general abelian-variety
bound has been asserted.

\paragraph{U06: boundaries of verification.}
The complete new proofs were read, and every included computational
certificate reports the exact execution actually performed. No
formal proof-assistant verification of this paper, reconstruction
of all earlier profinite elements, replay of every historical
computer search, or verification of Masullo's entire draft is
claimed. The check table distinguishes whole-proof reading,
independent arithmetic reimplementation, retained-verifier replay,
source-statement checks and carried background. A same-model second
or third reader is not presented as verification by a different
model or by a theorem prover.
